\scanpage{11}
\margin{10}

- 12 -\hfill n° 234

6. \uline{Extensions aux puissances tensorielles et extérieures.}

Supposons $A$ \uline{commutatif}, et soit $T^p(E)$ le $A$-module puissance tensorielle $p$-ème de $E$ (chap.III,\S\S1 et 4). Pour toute application semi-linéaire $u$ (relative à $J$) de $E$ dans un $A$-module $F$, il existe une application semi-linéaire et une seule $u'$ (relative à $J$) de $T^p(E)$ dans $T^p(F)$, telle que
\[
(18)\qquad u'(x_1\otimes x_2\otimes\ldots\otimes x_p)=u(x_1)\otimes u(x_2)\otimes\ldots\otimes u(x_p)
\]
Cela se voit en tenant compte de la définition du produit tensoriel, exactement comme pour le cas où $u$ est linéaire (si $v(x_1,\ldots,x_p)$ est le second-membre de (18), $v$ se prolonge en une application semi-linéaire $\bar v$ (relative à $J$) du module $G=A^{(E^p)}$ des combinaisons linéaires formelles des éléments de l'ensemble $E^p$, et comme
\[
v(x_1,\ldots,\lambda x_i,\ldots,x_p)=\lambda^J v(x_1,\ldots,x_p)
\]
$\bar v$ s'annule pour tous les éléments de $G$ de la forme
\[
(x_1,\ldots,\lambda x_i,\ldots,x_p)-\lambda(x_1,\ldots,x_p)\ ;
\]
cf.chap.III,\S1,n°2). Appliquons ce résultat aux applications semi linéaires $f$ et $g$ de $E$ dans $\unsure{E|E^*}$, associées à une forme sesquilinéaire $\Phi$ ; on obtient deux applications semi-linéaires de $T^p(E)$ dans $T^p(E^*)$ ; désignons par $f_p,g_p$ les applications semi-linéaires obtenues en composant avec ces applications, l'application canonique de $T^p(E^*)$ dans le dual $(T^p(E))^*$ (chap.III,\S1,n°4) ; on a par définition
\[
\begin{aligned}
\langle x_1\otimes\ldots\otimes x_p,f_p(y_1\otimes\ldots\otimes y_p)\rangle
&=\langle x_1,f(y_1)\rangle\ldots\langle x_p,f(y_p)\rangle\\
&=\Phi(x_1,y_1)\ldots\Phi(x_p,y_p)
\end{aligned}
\]
et de même
\[
\langle y_1\otimes\ldots\otimes y_p,g_p(x_1\otimes\ldots\otimes x_p)\rangle^J
=\langle y_1,g(x_1)\rangle^J\ldots\langle y_p,g(x_p)\rangle^J=
\]
